\section*{Capítulo \ref{ch:g4:weight-and-mass} --- El peso y la masa}

\begin{solution}{exo:g4:weight-and-mass:1}
La \omterm{def:g3:measuring-mass:mass}{masa}: ¿cuánta sustancia hay? El \omterm{def:g4:weight-and-mass:weight}{peso}: ¿con cuánta \omterm{def:g2:pushes-and-pulls:force}{fuerza} tiran
de ella aquí? El \omterm{def:g4:weight-and-mass:weight}{peso} es la \omterm{def:g2:pushes-and-pulls:force}{fuerza} --- el tirón de la Tierra
hacia abajo.
\end{solution}

\begin{solution}{exo:g4:weight-and-mass:2}
Hacia abajo, siempre. El \omterm{def:g4:weight-and-mass:springscale}{dinamómetro} --- el \omterm{def:g4:weight-and-mass:weight}{peso} estira su muelle
y la aguja informa del estiramiento.
\end{solution}

\begin{solution}{exo:g4:weight-and-mass:3}
Unos cuatro anchos de dedo --- el doble de sustancia, el doble de
tirón, el doble de estiramiento: la proposición ``más sustancia, tirón
más fuerte''.
\end{solution}

\begin{solution}{exo:g4:weight-and-mass:4}
El estiramiento hay que medirlo \emph{desde la posición en vacío}: la
goma ya cuelga un poco por el \omterm{def:g4:weight-and-mass:weight}{peso} del propio vaso, y solo el
descenso extra por debajo de la marca del cero mide la carga.
\end{solution}

\begin{solution}{exo:g4:weight-and-mass:5}
Sí --- la \omterm{def:g3:measuring-mass:mass}{masa} viaja sin cambiar: siguen siendo \qty{20}{kg} de
sustancia. Pero es muchísimo más fácil sostenerla: la \omterm{def:g2:moon-in-the-sky:moon}{Luna} tira de
ella unas seis veces más flojo, así que su \omterm{def:g4:weight-and-mass:weight}{peso} es alrededor de la
sexta parte del terrestre.
\end{solution}

\begin{solution}{exo:g4:weight-and-mass:6}
Sus piernas empujan tan fuerte como empujan siempre unas piernas
entrenadas en la Tierra, pero el tirón de vuelta de la \omterm{def:g2:moon-in-the-sky:moon}{Luna} es seis
veces más flojo --- así que cada paso corriente se convierte en un
salto.
\end{solution}

\begin{solution}{exo:g4:weight-and-mass:7}
El \omterm{def:g4:weight-and-mass:springscale}{dinamómetro}: nota el tirón local y marca seis veces menos. La
\omterm{def:g3:levers-and-scales:scale}{balanza} compara el tirón sobre la mochila con el tirón sobre las
pesas, y la \omterm{def:g2:moon-in-the-sky:moon}{Luna} debilita los dos por igual --- los platillos siguen
en tablas con la misma \omterm{def:g3:measuring-mass:mass}{masa}.
\end{solution}

\begin{solution}{exo:g4:weight-and-mass:8}
Verdadero. El \omterm{def:g4:weight-and-mass:weight}{peso} es el tirón de la Tierra sobre la sustancia de una
cosa, y el tirón crece con la cantidad de sustancia --- sin sustancia
no hay de qué tirar: \omterm{def:g4:weight-and-mass:weight}{peso} cero en todas partes.
\end{solution}

\begin{solution}{exo:g4:weight-and-mass:9}
El muelle de dentro nota el \emph{\omterm{def:g4:weight-and-mass:weight}{peso}} --- el tirón de las manzanas
sobre el platillo. En la Tierra, el tirón y la sustancia viajan
estrictamente juntos (más tirón exactamente cuando hay más
sustancia), así que una \omterm{def:g3:levers-and-scales:scale}{balanza} marcada en \omterm{def:g3:measuring-mass:units}{kilogramos} de sustancia
sigue diciendo la verdad --- aquí.
\end{solution}

\begin{solution}{exo:g4:weight-and-mass:10}
Las medidas de \omterm{def:g3:measuring-mass:mass}{masa} viajan sin cambiar: describen la sustancia misma,
no el sitio. Cualquier pesada hecha por \emph{comparación} --- la
\omterm{def:g3:levers-and-scales:scale}{balanza} contra pesas de referencia --- da la \omterm{def:g3:measuring-mass:mass}{masa} y sigue siendo
exacta cerca del planeta lejano. Solo se encogerán allí las lecturas
de \omterm{def:g4:weight-and-mass:weight}{peso} hechas con muelle, y no eran esas las que anotó la lista de
piezas.
\end{solution}

\begin{solution}{exo:g4:weight-and-mass:11}
La \omterm{def:g3:measuring-mass:mass}{masa} sigue toda ahí. Pistas: agitar la llave sigue costando
esfuerzo de verdad --- su mano nota la sustancia resistiéndose al
vaivén ---, y al lanzarla sale volando solo tan deprisa como sus
músculos consigan --- un martillo de sustancia sigue siendo un
martillo a la hora de lanzarlo. El \omterm{def:g4:weight-and-mass:weight}{peso} ha desaparecido; la
sustancia, jamás.
\end{solution}
